\documentclass[addpoints]{exam}

\usepackage{epic,array,ecltree,url,calrsfs}
\usepackage[nointegrals]{wasysym}

\usepackage{array,epsfig}
\usepackage{amsmath}
\usepackage{amsfonts}
\usepackage{amssymb}
\usepackage{amsxtra}
\usepackage{amsthm}
\usepackage{../mlextra} % must be below ams packages
\usepackage{mathrsfs}
\usepackage[dvipsnames]{xcolor}
\usepackage{array}
\usepackage{graphicx}
\graphicspath{ {../../art/} }
\usepackage{bm}
\usepackage{tikz}
\usepackage{multicol}

%Pagination stuff.
\setlength{\topmargin}{-.3 in}
\setlength{\oddsidemargin}{0in}
\setlength{\evensidemargin}{0in}
\setlength{\textheight}{9.in}
\setlength{\textwidth}{6.5in}

\newcommand{\tf}[1][{}]{%
\fillin[#1][0.25in]%
}

\noprintanswers
\unframedsolutions
\SolutionEmphasis{\itshape\small}
\SolutionEmphasis{\color{NavyBlue}}
\checkboxchar{$\Box$}
\checkedchar{$\blacksquare$}

%\pagestyle{empty} 
%\footer{}{\thepage}{}

\pointsinmargin
\begin{document}


\noindent
\begin{tabular*}{\textwidth}{l @{\extracolsep{\fill}} r @{\extracolsep{6pt}} l}
{\large CS3001: Foundations of Computer Science} &  \makebox[3in]{\large Name:\enspace\hrulefill}\\
{\large June 11, 2019} & \\
{\large Exam 3} & 
\end{tabular*}\\

\fbox{\fbox{\parbox{6in}{\textbf{Instructions}: Please answer the questions 
  below to the best of your ability. Be sure to show your work where appropriate. 
   This exam is closed book, closed notes, closed computer. There are \numpoints\ 
   points in total and \numbonuspoints\ possible extra credit points. }}}\\
\begin{questions}

%First section: 10 points
%
%Second section: 10 points
%
%Temporal Logic: 35 points
%
%Mutual exclusion property.
%Liveness Property.

%\question Define propositions $p$, $q$, $p'$ and $q'$ as follows:\\
% $p$: ``Process p is in the critical section.''\\
% $q$: ``Process q is in the critical section.''\\ 
% $p'$: ``Process p is attempting to enter the critical section.\\
% $q'$: ``Process q is attempting to enter the critical section.''\\
%\begin{parts}
%\part[2] Use the propositions given above to express the property of
%mutual exclusion in linear temporal logic.
%\vspace{15mm}
%
%\begin{solution}
%$\Box \ngg (p \land q)$
%\end{solution}
%
%\part[2] Use the propositions given above to express the property of
%liveness in linear temporal logic.
%\vspace{15mm}
%
%\begin{solution}
%$\Box (p' \imp \Diamond p) \land \Box (q' \imp \Diamond q)$
%\end{solution}
%\end{parts}
%
\question For each formula below select all that apply:

\begin{parts}
\part[2] $\Circle p \land \ngg \Diamond  p$

\begin{oneparcheckboxes}
\choice Valid
\choice Satisfiable
\CorrectChoice Falsifiable
\CorrectChoice Unsatisfiable
\end{oneparcheckboxes}
\vspace{5mm}

\part[2] $\Box(p \imp \Circle p) \imp \Box p$

\begin{oneparcheckboxes}
\choice Valid
\CorrectChoice Satisfiable
\CorrectChoice Falsifiable
\choice Unsatisfiable
\end{oneparcheckboxes}
\vspace{5mm}

\part[2] $\Box(p \land q) \eqv (\Box p \land \Box q)$

\begin{oneparcheckboxes}
\CorrectChoice Valid
\CorrectChoice Satisfiable
\choice Falsifiable
\choice Unsatisfiable
\end{oneparcheckboxes}
\vspace{5mm}

\part[2] $\Diamond \Box \ngg p \imp \Box \Diamond p $   

\begin{oneparcheckboxes}
\choice Valid
\choice Satisfiable
\CorrectChoice Falsifiable
\CorrectChoice Unsatisfiable
\end{oneparcheckboxes}
\vspace{5mm}

\end{parts}
\question[4] Consider the following interpretation $\mathcal{I} = (P,\{C,E\},\{Fred\})$,
where 
\begin{itemize}
\item $P$ is a set of people, 
\item $C$ is a unary relation such that $C(x)$ iff $x$ is a civilian, 
\item $E$ is a unary relation such that $E(x)$ iff $x$ may enter the library.
\item Fred is some person.
\end{itemize}
Write a statement in first order logic to express the following:
\begin{parts}
\part All civilians may enter the library if only civilians may enter the
library.
\vspace{15mm}
\part If Fred is a civilian, then no one may enter the library.
\vspace{15mm}
\end{parts}
\question[3] Let $\Sigma = \{a\}$ and let $K$ be the set of odd length
strings over $\Sigma^*$. Write a Post system for $K$.

\clearpage
\question For each of the formulas and sets of interpretations below, mark $S$
next to each interpretation that {\bf \emph{satisfies}} the formula, and $F$ next to each
interpretation that {\bf \emph{falsifies}} the formula. (For all interpretations shown,
assume the values of the last two states continue to repeat after $s_4$, so 
$s_5 = s_3$, $s_6 = s_4$, etc.)
\begin{parts}
\part[4] $\Box (p \lor q) \imp (\Box p \lor \Box q)$ \\
\begin{picture}(340,1)
\put(200, 0){\makebox(25,20)[l]{\small Repeat $\imp$}}
\multiput(190, -135)(0,5){30}{\line(0,1){1}}
\end{picture}
\begin{choices}
\choice
\unitlength=.9pt
\begin{picture}(340,30)
\put(  0,-5){\state{$p$}{$s_{0}$}}
\put( 20, 5){\vector(1,0){40}}
\put( 60,-5){\state{$p$}{$s_{1}$}}
\put( 80, 5){\vector(1,0){40}}
\put(120,-5){\state{$p$}{$s_{2}$}}
\put(140, 5){\vector(1,0){40}}
\put(180,-5){\state{$p$}{$s_{3}$}}
\put(200, 5){\vector(1,0){40}}
\put(240,-5){\state{$p$}{$s_{4}$}}
\put(260, 5){\vector(1,0){40}}
\put(300, 0){\makebox(20,10){\ldots}}
\end{picture}
\tf[S] (S/F)

\choice
\unitlength=.9pt
\begin{picture}(340,30)
\put(  0,-5){\state{$p$}{$s_{0}$}}
\put( 20, 5){\vector(1,0){40}}
\put( 60,-5){\state{$q$}{$s_{1}$}}
\put( 80, 5){\vector(1,0){40}}
\put(120,-5){\state{$p$}{$s_{2}$}}
\put(140, 5){\vector(1,0){40}}
\put(180,-5){\state{$q$}{$s_{3}$}}
\put(200, 5){\vector(1,0){40}}
\put(240,-5){\state{$p$}{$s_{4}$}}
\put(260, 5){\vector(1,0){40}}
\put(300, 0){\makebox(20,10){\ldots}}
\end{picture}
\tf[F] (S/F)



\choice
\unitlength=.9pt
\begin{picture}(340,30)
\put(  0,-5){\state{}{$s_{0}$}}
\put( 20, 5){\vector(1,0){40}}
\put( 60,-5){\state{$q$}{$s_{1}$}}
\put( 80, 5){\vector(1,0){40}}
\put(120,-5){\state{$p$}{$s_{2}$}}
\put(140, 5){\vector(1,0){40}}
\put(180,-5){\state{$q$}{$s_{3}$}}
\put(200, 5){\vector(1,0){40}}
\put(240,-5){\state{$p$}{$s_{4}$}}
\put(260, 5){\vector(1,0){40}}
\put(300, 0){\makebox(20,10){\ldots}}
\end{picture}
\tf[S] (S/F)

\choice 
\unitlength=.9pt
\begin{picture}(340,30)
\put(  0,-5){\state{}{$s_{0}$}}
\put( 20, 5){\vector(1,0){40}}
\put( 60,-5){\state{}{$s_{1}$}}
\put( 80, 5){\vector(1,0){40}}
\put(120,-5){\state{}{$s_{2}$}}
\put(140, 5){\vector(1,0){40}}
\put(180,-5){\state{}{$s_{3}$}}
\put(200, 5){\vector(1,0){40}}
\put(240,-5){\state{}{$s_{4}$}}
\put(260, 5){\vector(1,0){40}}
\put(300, 0){\makebox(20,10){\ldots}}
\end{picture}
\tf[S] (S/F)


\end{choices}


\part[4] $(\Box p \imp \Box q) \imp \Box( p \imp q) $\\
\begin{picture}(340,1)
\put(200, 0){\makebox(25,20)[l]{\small Repeat $\imp$}}
\multiput(190, -135)(0,5){30}{\line(0,1){1}}
\end{picture}
\begin{choices}

\choice 
\unitlength=.9pt
\begin{picture}(340,30)
\put(  0,-5){\state{$pq$}{$s_{0}$}}
\put( 20, 5){\vector(1,0){40}}
\put( 60,-5){\state{$pq$}{$s_{1}$}}
\put( 80, 5){\vector(1,0){40}}
\put(120,-5){\state{$pq$}{$s_{2}$}}
\put(140, 5){\vector(1,0){40}}
\put(180,-5){\state{$pq$}{$s_{3}$}}
\put(200, 5){\vector(1,0){40}}
\put(240,-5){\state{$pq$}{$s_{4}$}}
\put(260, 5){\vector(1,0){40}}
\put(300, 0){\makebox(20,10){\ldots}}
\end{picture}
\tf[S] (S/F)

\choice
\unitlength=.9pt
\begin{picture}(340,30)
\put(  0,-5){\state{$q$}{$s_{0}$}}
\put( 20, 5){\vector(1,0){40}}
\put( 60,-5){\state{$q$}{$s_{1}$}}
\put( 80, 5){\vector(1,0){40}}
\put(120,-5){\state{$p$}{$s_{2}$}}
\put(140, 5){\vector(1,0){40}}
\put(180,-5){\state{$q$}{$s_{3}$}}
\put(200, 5){\vector(1,0){40}}
\put(240,-5){\state{$q$}{$s_{4}$}}
\put(260, 5){\vector(1,0){40}}
\put(300, 0){\makebox(20,10){\ldots}}
\end{picture}
\tf[F] (S/F)


\choice
\unitlength=.9pt
\begin{picture}(340,30)
\put(  0,-5){\state{$q$}{$s_{0}$}}
\put( 20, 5){\vector(1,0){40}}
\put( 60,-5){\state{$pq$}{$s_{1}$}}
\put( 80, 5){\vector(1,0){40}}
\put(120,-5){\state{}{$s_{2}$}}
\put(140, 5){\vector(1,0){40}}
\put(180,-5){\state{$q$}{$s_{3}$}}
\put(200, 5){\vector(1,0){40}}
\put(240,-5){\state{$pq$}{$s_{4}$}}
\put(260, 5){\vector(1,0){40}}
\put(300, 0){\makebox(20,10){\ldots}}
\end{picture}
\tf[F] (S/F)

\choice
\unitlength=.9pt
\begin{picture}(340,30)
\put(  0,-5){\state{$p$}{$s_{0}$}}
\put( 20, 5){\vector(1,0){40}}
\put( 60,-5){\state{$p$}{$s_{1}$}}
\put( 80, 5){\vector(1,0){40}}
\put(120,-5){\state{$p$}{$s_{2}$}}
\put(140, 5){\vector(1,0){40}}
\put(180,-5){\state{$p$}{$s_{3}$}}
\put(200, 5){\vector(1,0){40}}
\put(240,-5){\state{$p$}{$s_{4}$}}
\put(260, 5){\vector(1,0){40}}
\put(300, 0){\makebox(20,10){\ldots}}
\end{picture}
\tf[F] (S/F)
\end{choices}

\part[4] $\Box (q \imp \Box \Diamond p)$ \\
%\part[4] $\Box \Diamond p \imp \Diamond \Box p$ \\
\begin{picture}(340,1)
\put(200, 0){\makebox(25,20)[l]{\small Repeat $\imp$}}
\multiput(190, -135)(0,5){30}{\line(0,1){1}}
\end{picture}
\begin{choices}
\choice
\unitlength=.9pt
\begin{picture}(340,30)
\put(  0,-5){\state{$q$}{$s_{0}$}}
\put( 20, 5){\vector(1,0){40}}
\put( 60,-5){\state{}{$s_{1}$}}
\put( 80, 5){\vector(1,0){40}}
\put(120,-5){\state{$p$}{$s_{2}$}}
\put(140, 5){\vector(1,0){40}}
\put(180,-5){\state{$q$}{$s_{3}$}}
\put(200, 5){\vector(1,0){40}}
\put(240,-5){\state{}{$s_{4}$}}
\put(260, 5){\vector(1,0){40}}
\put(300, 0){\makebox(20,10){\ldots}}
\end{picture}
\tf[F] (S/F)


\choice 
\unitlength=.9pt
\begin{picture}(340,30)
\put(  0,-5){\state{$q$}{$s_{0}$}}
\put( 20, 5){\vector(1,0){40}}
\put( 60,-5){\state{}{$s_{1}$}}
\put( 80, 5){\vector(1,0){40}}
\put(120,-5){\state{}{$s_{2}$}}
\put(140, 5){\vector(1,0){40}}
\put(180,-5){\state{}{$s_{3}$}}
\put(200, 5){\vector(1,0){40}}
\put(240,-5){\state{}{$s_{4}$}}
\put(260, 5){\vector(1,0){40}}
\put(300, 0){\makebox(20,10){\ldots}}
\end{picture}
\tf[F] (S/F)

\choice
\unitlength=.9pt
\begin{picture}(340,30)
\put(  0,-5){\state{$p$}{$s_{0}$}}
\put( 20, 5){\vector(1,0){40}}
\put( 60,-5){\state{}{$s_{1}$}}
\put( 80, 5){\vector(1,0){40}}
\put(120,-5){\state{}{$s_{2}$}}
\put(140, 5){\vector(1,0){40}}
\put(180,-5){\state{$p$}{$s_{3}$}}
\put(200, 5){\vector(1,0){40}}
\put(240,-5){\state{$p$}{$s_{4}$}}
\put(260, 5){\vector(1,0){40}}
\put(300, 0){\makebox(20,10){\ldots}}
\end{picture}
\tf[S] (S/F)
  

\choice
\unitlength=.9pt
\begin{picture}(340,30)
\put(  0,-5){\state{$p$}{$s_{0}$}}
\put( 20, 5){\vector(1,0){40}}
\put( 60,-5){\state{}{$s_{1}$}}
\put( 80, 5){\vector(1,0){40}}
\put(120,-5){\state{}{$s_{2}$}}
\put(140, 5){\vector(1,0){40}}
\put(180,-5){\state{$p$}{$s_{3}$}}
\put(200, 5){\vector(1,0){40}}
\put(240,-5){\state{$q$}{$s_{4}$}}
\put(260, 5){\vector(1,0){40}}
\put(300, 0){\makebox(20,10){\ldots}}
\end{picture}
\tf[S] (S/F)



\end{choices}
\part[4] $(\Diamond \Box p \land \Box \Diamond q) \imp \Box \Diamond (p \land q)$\\
\begin{picture}(340,1)
\put(200, 0){\makebox(25,20)[l]{\small Repeat $\imp$}}
\multiput(190, -135)(0,5){30}{\line(0,1){1}}
\end{picture}

\begin{choices}

\choice 
\unitlength=.9pt
\begin{picture}(340,30)
\put(  0,-5){\state{$p$}{$s_{0}$}}
\put( 20, 5){\vector(1,0){40}}
\put( 60,-5){\state{}{$s_{1}$}}
\put( 80, 5){\vector(1,0){40}}
\put(120,-5){\state{$p$}{$s_{2}$}}
\put(140, 5){\vector(1,0){40}}
\put(180,-5){\state{$p$}{$s_{3}$}}
\put(200, 5){\vector(1,0){40}}
\put(240,-5){\state{$pq$}{$s_{4}$}}
\put(260, 5){\vector(1,0){40}}
\put(300, 0){\makebox(20,10){\ldots}}
\end{picture}
\tf[S] (S/F)


\choice
\unitlength=.9pt
\begin{picture}(340,30)
\put(  0,-5){\state{$p$}{$s_{0}$}}
\put( 20, 5){\vector(1,0){40}}
\put( 60,-5){\state{}{$s_{1}$}}
\put( 80, 5){\vector(1,0){40}}
\put(120,-5){\state{}{$s_{2}$}}
\put(140, 5){\vector(1,0){40}}
\put(180,-5){\state{$q$}{$s_{3}$}}
\put(200, 5){\vector(1,0){40}}
\put(240,-5){\state{$p$}{$s_{4}$}}
\put(260, 5){\vector(1,0){40}}
\put(300, 0){\makebox(20,10){\ldots}}
\end{picture}
\tf[S] (S/F)

\choice
\unitlength=.9pt
\begin{picture}(340,30)
\put(  0,-5){\state{$pq$}{$s_{0}$}}
\put( 20, 5){\vector(1,0){40}}
\put( 60,-5){\state{$pq$}{$s_{1}$}}
\put( 80, 5){\vector(1,0){40}}
\put(120,-5){\state{$pq$}{$s_{2}$}}
\put(140, 5){\vector(1,0){40}}
\put(180,-5){\state{$p$}{$s_{3}$}}
\put(200, 5){\vector(1,0){40}}
\put(240,-5){\state{$p$}{$s_{4}$}}
\put(260, 5){\vector(1,0){40}}
\put(300, 0){\makebox(20,10){\ldots}}
\end{picture}
\tf[S] (S/F)



\choice
\unitlength=.9pt
\begin{picture}(340,30)
\put(  0,-5){\state{}{$s_{0}$}}
\put( 20, 5){\vector(1,0){40}}
\put( 60,-5){\state{}{$s_{1}$}}
\put( 80, 5){\vector(1,0){40}}
\put(120,-5){\state{}{$s_{2}$}}
\put(140, 5){\vector(1,0){40}}
\put(180,-5){\state{}{$s_{3}$}}
\put(200, 5){\vector(1,0){40}}
\put(240,-5){\state{}{$s_{4}$}}
\put(260, 5){\vector(1,0){40}}
\put(300, 0){\makebox(20,10){\ldots}}
\end{picture}
\tf[S] (S/F)

\end{choices}



\end{parts}

\clearpage
%\question[5] Consider the following attempt to solve the critical section
%problem:
%\begin{center}
%\begin{tabular}{|p{0.45\textwidth}|p{0.45\textwidth}|}
%\hline
%\multicolumn{2}{|c|}{\p{boolean wantp = false, wantq = false}}\\
%\hline
%\p{Process p} & \p{Process q} \\
%\hline
%\p{while (true) \{} & \p{while (true) \{} \\
%\p{\ tryp: \ wantp = true} & \p{\ tryq: \ wantq = true } \\
%\p{\ waitp: wait until !wantq } & \p{\ waitq: wait until !wantp} \\
%\p{\ csp: \ \ wantp = false} & \p{\ csq: \ \ wantq = false} \\
%\p{\}} & \p{\}} \\\hline
%\end{tabular}
%\end{center}
%Draw the state transition diagram for process $p$.
%\vspace{30mm}
\question[4] Let $A$ be some set of formulas in first order logic and $B$ the
set of all $\id{FOL}$ formulas that are not in $A$. Which of
the following can you conclude (choose all that apply)?
\begin{checkboxes}
\CorrectChoice If $A$ is decidable, $B$ is decidable.
\choice If $A$ is not decidable, $B$ is semi-decidable.
\choice If $A$ is in $P$, $B$ is semi-decidable.
\CorrectChoice If $A$ is semi-decidable, $B$ is neither decidable nor semi-decidable.
\end{checkboxes}
\vspace{5mm}

\question The following is proposed as an alternative one-bit adder design:

\begin{center}
\begin{picture}(280,120)
\put(-10,  0){\makebox(20,20)[l]{$b_2$}}
\put(-10,100){\makebox(20,20)[l]{$b_1$}}
\put(260, 20){\makebox(25,20)[l]{s}}
\put(260, 80){\makebox(25,20)[l]{c}}
\put(20, 10){\line(1,0){50}}
\put(20,110){\line(1,0){60}}
\put(30, 10){\line(1, 2){40}}
\put(30,110){\line(1,-2){40}}
\put(70, 90){\line(1,0){10}}
\put(70,80){ %OR gate
  \put(20,20){\oval(40,40)[r]}
  \put(0,20){\oval(20,40)[r]}
  \put(0,0){\line(1,0){20}}
  \put(0,40){\line(1,0){20}}
}
\put(70,0){ %AND gate
  \put(20,20){\oval(40,40)[r]}
  \put(0,0){\line(0,1){40}}
  \put(0,0){\line(1,0){20}}
  \put(0,40){\line(1,0){20}}
}
\put(110,100){\line(1,0){60}}
\put(110, 20){\line(1,0){60}}
\put(210, 30){\line(1,0){40}}
\put(210, 90){\line(1,0){40}}
\put(170,70){ %AND gate
  \put(20,20){\oval(40,40)[r]}
  \put(0,0){\line(0,1){40}}
  \put(0,0){\line(1,0){20}}
  \put(0,40){\line(1,0){20}}
}\put(170,10){ %AND gate
  \put(20,20){\oval(40,40)[r]}
  \put(0,0){\line(0,1){40}}
  \put(0,0){\line(1,0){20}}
  \put(0,40){\line(1,0){20}}
}
\put(120,80){  % NOT gate
  \put( 0, 0){\line(1,0){20}}
  \put( 0, 0){\line(1,-2){10}}
  \put(20, 0){\line(-1,-2){10}}
  \put(10,-23){\circle{6}}
}
\put(130,80){\line(0,1){20}}
\put(130,54){\line(0,-1){14}}
\put(130,40){\line(1,0){40}}
\put(150,80){\line(1,0){20}}
\put(150,20){\line(0,1){60}}
\end{picture}
\end{center}

\begin{parts}
\part[4] Write a formula in propositional logic that relates: 

\begin{subparts}

\subpart $b_1$, $b_2$ and $c$
\vspace{15mm}
\begin{solution}
$(b_1 \lor b_2) \land (b_1 \land b_2) \eqv c$
\end{solution}

\subpart $b_1$, $b_2$ and $s$
\vspace{15mm}
\begin{solution}
$\ngg (b_1 \lor b_2) \land (b_1 \land b_2) \eqv s$
\end{solution}
\end{subparts}

\part[4] Assume $s$ and $c$ are meant to model the behavior of sum and carry
bits respectively, and use your statements to evaluate the proposed design.

\begin{subparts}
\subpart Will the sum bit function correctly? If yes, explain how you know
this. If no, provide an example.
\vspace{15mm}
\subpart Will the carry bit function correctly? If yes, explain how you know
this. If no, provide an example.
\vspace{15mm}
\end{subparts}


\end{parts}


\clearpage
\question Answer parts a, b and c with respect to the automaton shown below:
\fullwidth{\begin{center}
\includegraphics{exam3_automaton.pdf}
\end{center}}

\begin{parts}
\part[4] Fill in the blanks to complete the definition of the transition
function, $\delta$.

\begin{align*}
  &\delta(s_0,0) = s_0&
  &\delta(s_1,\fillin[$0$][.25in]) = \fillin[$s_2$][.25in]
  &\delta(s_2,\fillin[$0$][.25in]) = \fillin[$s_2$][.25in]\\ 
  &\delta(s_0,1) = s_1&
  &\delta(s_1,\fillin[$1$][.25in]) = \fillin[$s_2$][.25in]
  &\delta(s_2,\fillin[$1$][.25in]) = \fillin[$s_2$][.25in]\\
\end{align*}

\part[2] Let $\hat{\delta}(s_1,w) = s_0$. Which of the following are possible
values of $w$ (circle all that apply)?

\begin{oneparchoices}
\CorrectChoice $11$
\choice $11110$
\CorrectChoice $1111011$
\choice $100000111$
\end{oneparchoices}
\vspace{5mm}

\part[3] Say we want to prove that this automaton accepts only strings
ending in an even number of ones. Our strategy will be to perform
induction on input length. Fill in the blanks below to complete the statements
we will establish to prove the theorem:
\begin{itemize}
\item The automaton halts in state $s_0$ iff its input ends with a \fillin[$0$][.25in]. 
\item The automaton halts in state $s_1$ iff its input ends with an \fillin[odd][.25in] number of ones.
\item The automaton halts in state $s_2$ iff \fillin[its input ends with an odd
number of ones] [3in].
\end{itemize}
\part[3] Which of the following induction rules is the best choice for our proof?
\begin{choices}
\uplevel{
\choice 
	\(\begin{array}[t]{l}
	S(i) \;\wedge\; \forall n \geq i.\;S(n) \imp S(n+1) \\
	\hline
	\forall n \geq i. \; S(n)
	\end{array}\) % \\[2ex]

\CorrectChoice 
	\(\begin{array}[t]{l}
	S_1(i)\;\wedge\;S_2(i)\;\wedge\;S_3(i)\;\wedge\;\forall n \geq
        i.\;S_1(n)\;\wedge\;S_2(n)\;\wedge\;S_3(n)\imp S_1(n+1)\;\wedge\;S_2(n+1)\wedge\;S_3(n+1)\\
	\hline
	\forall n \geq i. \; S_1(n)\;\wedge\;S_2(n)\;\wedge\;S_3(n)

	\end{array}\) % \\[2ex]

\choice 
	\(\begin{array}[t]{l}
	S_1(i) \;\wedge\;S_2(i)\;\wedge\;S_3(i)  \;\wedge\; \\
        \forall n \geq i.\;S_1(i)\;\wedge\;S_2(i)\;\wedge\;S_3(i)\;\wedge\cdots\wedge\;S_1(n)\;\wedge\;S_2(n)\;\wedge\;S_3(n)\imp S_1(n+1)\;\wedge\;S_2(n+1)\;\wedge\;S_3(n+1)\\
	\hline
	\forall n \geq i. \; S_1(n)\;\wedge\;S_2(n)\;\wedge\;S_3(n)
	\end{array}\)

\choice 
	\(\begin{array}[t]{l}
	S_1(i) \;\wedge\; S_2(i)\;\wedge\;S_1(i+1) \;\wedge\; S_2(i+1)\;\wedge\;\cdots\wedge\;S_1(j)\;\wedge\;S_2(j) \\
\forall n \geq
j.\;S_1(i)\;\wedge\;S_2(i)\;\wedge\cdots\wedge\;S_1(n)\wedge\;S_2(n)\imp S_1(n+1)\wedge\;S_2(n+1)\\
	\hline
	\forall n \geq i. \; S_1(n)\;\wedge\;S_2(n)
	\end{array}\) % \\[2ex]

}
\end{choices}

\part[2] Briefly justify your answer to part (c):

\end{parts}

\clearpage
\question[5] Consider the following set of formulas in first order logic:
\begin{itemize}
\item $\forall x p(x,x)$
\item $\forall x \forall y (p(x,y) \land p(y,x) \imp q(x,y))$
\item $\forall x \forall y \forall z ((p(x,y) \land p(y,z)) \imp p(x,z))$
\end{itemize}
Which of the following represent \textbf{satisfying} interpretations?
Choose all that apply:
\begin{choices} 
\choice$\mathcal{I}_A = \{\Z, \{ \leq , = \}, \{\} \}$ 
\choice$\mathcal{I}_A = \{\N, \{ > , = \}, \{\} \}$ 
\choice$\mathcal{I}_A = \{\N, \{ \;|\; , = \}, \{\} \}$ (Note that $|$ indicates ``divides'')
\choice$\mathcal{I}_A = \{\{\emptyset, \{1\}, \{1,2\}\}, \{\subseteq , =  \}, \{\} \}$ 
\choice$\mathcal{I}_A = \{\mathcal{F}, \{\;\models \;, \;\equiv \; \}, \{\} \}$
(Let $\mathcal{F}$ be the set of formulas in propositional logic.)
\end{choices}

%\question[4] Let $S_{rt}$ be the set of reflexive, transitive relations, and let
%$\rho^*$ be a relation on the set of boolean functions $F_B$ such that 
%$\rho^* \in S_{rt}$.  Use first order logic to express the properties of $\rho^*$. 
\begin{solution}
$(\forall x (S(x) \imp \rho^*(x,x))) \land (\forall x\forall y \forall z (S(x)
      \land S(y) \land S(z) \land \rho^*(x,y) \land rho^*(y,z) \imp \rho^*(x,z)))$
\end{solution}
\vspace{5mm}


\bonusquestion[4] 
Let $\Sigma = \{a,b\}$ and let $p(x,y)$ be a binary relation on $\Sigma^*$ 
where $p(x,y)$ if and only if there exists a string $w$ in $\Sigma^*$ such that $y = xw$.
For example, if $x = aba$ and $y = ababb$, then $p(x,y)$ because $w = bb$
and $bb \in \Sigma^*$.  Prove that $P$ is a partial order.
\vspace{30mm}


%\bonusquestion[4] Prove: $\models \Box (p \imp q) \imp (\Box p \imp \Box q)$


%Draw derivation tree / post system?
%Complexity questions?


%For final, need to add interpretations
%\question[4] Which of the following statements in first order logic correctly
%             expresses the anti-symmetric property (choose all that apply):
%\begin{checkboxes}
%\CorrectChoice $\forall x \forall y(p(x,y) \land p(y,x) \imp q(x,y)$
%\choice        $\forall x \forall y (\ngg q(x,y) \imp (p(x,y) \imp \ngg %p(y,x)))
%\end{checkboxes}

%\question[6] Let $S$ be a propositional logic formula in clausal form, such that
%$S = \{s\bar{t}, p\bar{r}\bar{s}, \bar{p}\bar{r}, r\bar{s}, \bar{p}rs, st, pt,
%  rt, t\}$. Use DPLL to prove $S$ is a contradiction. Draw the execution tree.
%
%\vspace{60mm}
%
%
%
%
%
%\question[4] Let $A$ and $B$ be problems such that $A$ reduces to $B$, and
%evaluate each of the following claims:
%\begin{parts}
%\part \tf[T] (T/F) If $B$ is decidable, $A$ is decidable.
%\part \tf[F] (T/F) If $B$ is NP-Hard, $A$ is decidable.
%\part \tf[T] (T/F) If FOL-VAL reduces to $A$, $B$ is undecidable.
%\part \tf[F] (T/F) If $B$ reduces to FOL-SAT, $B$ is undecidable.
%\end{parts}
%
%\vspace{5mm}
%
%\question Let $\mathcal{F}$ be the set of all formulas in propositional logic.
%Define $A \subseteq \mathcal{F}$ to be the set of formulas that your instructor
%finds interesting, and let $B$ be the set of all other formulas.
%\begin{parts}
%
%\part[4] Assume we show that a non-deterministic Turing machine can determine
%whether a formula is in $A$ in polynomial time. Which of the following can you 
%conclude (choose all that apply)?
%\begin{checkboxes}
%\CorrectChoice $A$ is in NP. 
%\choice        $A$ is NP-Hard.
%\CorrectChoice $A$ is decidable. 
%\CorrectChoice $A$ is semi-decidable. 
%\end{checkboxes}
%
%\vspace{5mm}
%\part[4] Now assume PL-SAT reduces to $A$ in polynomial time. Which of the following can you 
%conclude (choose all that apply)?
%\begin{checkboxes}
%\CorrectChoice $A$ is NP-Complete.
%\CorrectChoice $A$ is NP-Hard.
%\CorrectChoice $B$ is in CoNP.
%\CorrectChoice $B$ is CoNP-Complete.
%\end{checkboxes}
%\end{parts}
%
%\vspace{5mm}
%\question Let $FOL$ be the set of all formulas in first-order logic, and $A$-FOL  



%\part[1] $A = (\forall x p(x) \imp \forall x q(x)) \imp \forall x (p(x) \imp q(x))$ and $\mathcal{I}_A = \{\N_0, \{=\}, \{\times\}, \fillin \}$

\end{questions}
\end{document}


